\documentclass[11pt]{article}

\usepackage[margin=1in]{geometry}
\usepackage[T1]{fontenc}
\usepackage[utf8]{inputenc}
\usepackage{lmodern}
\usepackage{microtype}
\usepackage{hyperref}
\usepackage{natbib}
\usepackage{enumitem}

\setlist[itemize]{leftmargin=*, itemsep=0.25em}
\setlist[enumerate]{leftmargin=*, itemsep=0.25em}

\hypersetup{
  colorlinks=true,
  linkcolor=blue,
  urlcolor=blue,
  citecolor=blue,
  pdftitle={Why Safety-First AI Governance Risks Producing Unsafe Systems},
  pdfauthor={Anthony Paterson}
}

\title{Why Safety-First AI Governance Risks Producing Unsafe Systems\\\large Agency Suppression, Oversight Illusions, and Institutional Fragility}
\author{Anthony Paterson\\Fractal Media Infrastructure (FMI)\\\url{https://github.com/instance001}}
\date{2026-02-10}

\begin{document}
\maketitle

\noindent\textbf{Companion framework:} \texttt{asewb/} (Agency Suppression Early Warning Benchmark).

\begin{abstract}
Contemporary artificial intelligence governance is increasingly framed around
safety-first principles emphasizing restriction, limitation, and centralized control.
While such approaches aim to prevent misuse and harm, this paper argues that an
overreliance on safety-centric constraint risks reproducing a broader institutional
pattern: the suppression of human agency in the name of risk mitigation.

Drawing on a structural analysis of safety theatre, the paper examines how AI
governance frameworks may substitute procedural restriction for competence development,
creating systems that appear controlled while becoming brittle under real-world
complexity. We argue that safety in AI systems depends not only on limiting machine
behavior, but on preserving and cultivating human judgment, interpretability, and
intervention capacity. The paper concludes that agency is not an optional add-on to AI
safety, but a foundational requirement for resilient oversight.
\end{abstract}

\noindent\textbf{Keywords:} AI governance; AI safety; risk management; agency; oversight; institutional design

\section{Introduction: The Safety Reflex in AI Governance}

Artificial intelligence has emerged into public consciousness accompanied by an
unusually intense focus on safety. From the earliest large-scale deployments, discourse
around AI has emphasized alignment, guardrails, risk prevention, and misuse avoidance.
This emphasis is understandable: AI systems operate at scale, act with speed, and can
produce non-obvious harms.

In parallel, governance approaches increasingly emphasize risk management and control
layers, with an understandable focus on compliance and measurable assurance (e.g.,
\citet{NIST2023AIRMF}). This paper is not a critique of AI safety as such; for technical
problem framings within AI safety research, see \citet{Amodei2016Concrete}. The question
is not whether AI should be governed, but how safety is conceptualized---and what is
displaced in the process.

The analysis is framed as a structural one: how institutional safety narratives can
shift from harm reduction toward discretion control. In this sense, it is continuous
with work on legibility and local knowledge \citep{Scott1998Seeing} and with critiques
of risk-management and audit cultures \citep{Power2004Risk}.

\section{From Harm Prevention to Agency Restriction}

Many AI safety regimes implicitly treat human agency as a risk surface. Human
discretion, experimentation, and interpretation are framed as sources of
unpredictability to be minimized. As a result, governance mechanisms often prioritize:

\begin{itemize}
  \item Centralized approval processes
  \item Rigid usage constraints
  \item Oversight layers removed from operational context
  \item Procedural compliance over outcome evaluation
\end{itemize}

These mechanisms can reduce visible misuse but also reduce opportunities for humans to
develop operational understanding of AI systems. Over time, safety becomes defined not
by effective human--system interaction, but by the absence of deviation.

\section{Oversight Illusions and the Deskilling of Humans}

A key risk in safety-first AI governance is the creation of \emph{oversight illusions}.
These arise when systems appear well-governed because controls exist, even as
meaningful human understanding and intervention capacity erode.

Common examples include:

\begin{itemize}
  \item ``Human-in-the-loop'' roles that are symbolic rather than substantive
  \item Audit mechanisms focused on documentation rather than comprehension
  \item Safety reviews conducted by actors insulated from real-world deployment contexts
\end{itemize}

As discretion is removed, humans lose opportunities to practice judgment under
uncertainty. This produces a deskilling effect: people become less capable of
recognizing failure modes, interpreting anomalous behavior, or intervening effectively
when systems behave unexpectedly.

\section{Competence Suppression in AI Contexts}

Restriction can be appropriate, but restriction without parallel competence cultivation
produces a familiar loop:

\begin{enumerate}
  \item Interaction is limited to prevent error.
  \item Limited interaction reduces understanding.
  \item Reduced understanding increases perceived risk.
  \item Increased perceived risk justifies further restriction.
\end{enumerate}

This loop results in AI systems that are heavily constrained yet poorly understood by
those responsible for them. Safety becomes dependent on abstraction layers rather than
lived expertise. Such systems function adequately under expected conditions but fail
disproportionately under novelty---where predefined safeguards are least effective.

\section{Brittleness, Not Alignment, as a Dominant Failure Mode}

Many AI safety debates focus on alignment failures: systems behaving in ways misaligned
with human values. While alignment remains a real concern, brittleness may be a more
immediate institutional risk.

Brittle systems are:

\begin{itemize}
  \item Overfit to known scenarios
  \item Dependent on procedural compliance
  \item Incapable of graceful degradation
\end{itemize}

When AI governance suppresses human agency, it removes the adaptive buffer that absorbs
uncertainty. In such environments, failures propagate rapidly because no actor is
empowered---or practiced---to respond contextually.

This concern is closely related to classic observations about complex, tightly coupled
systems: controls that improve routine performance can still leave organizations unable
to recover under surprise \citep{Perrow1984Normal}.

\section{Agency as a Safety Requirement in AI Systems}

Agency in AI governance does not imply unrestricted access or unbounded experimentation.
Rather, it refers to the presence of humans who are:

\begin{itemize}
  \item Permitted to form internal models of system behavior
  \item Trained through interaction, not merely instruction
  \item Empowered to intervene meaningfully
  \item Accountable for outcomes, not just compliance
\end{itemize}

Safety emerges from disciplined agency, not from its absence. AI systems embedded in
organizations that discourage judgment may appear safe while accumulating hidden
fragility.

\section{Implications for AI Policy and System Design}

Reframing agency as a safety primitive suggests practical design implications:

\begin{itemize}
  \item Balance restriction with structured exposure and learning.
  \item Prioritize understanding over documentation in oversight.
  \item Ensure human roles are substantive, not symbolic.
  \item Design for recovery and adaptation, not only prevention.
\end{itemize}

Without these elements, AI governance risks becoming a specialized instance of safety
theatre---effective at limiting action, ineffective at handling reality.

\section{Conclusion: Unsafe by Design}

Safety-first AI governance often assumes that fewer decisions mean fewer mistakes. This
paper argues the opposite: systems that suppress human agency reduce their capacity to
detect, interpret, and respond to failure.

AI systems do not operate in isolation. Their safety depends on the humans who design,
deploy, and oversee them. When those humans are trained to defer rather than judge,
safety becomes performative rather than real.

In AI, as elsewhere, safety without agency is not safety. It is delay.

\section{Observable indicators of agency-suppressive safety drift}

Analyses of institutional safety drift often rely on retrospective interpretation. To
mitigate this limitation, this section outlines a set of observable interactional
indicators that may signal the emergence of safety theatre within AI systems and their
governance environments. These indicators are presented as probabilistic markers rather
than evidence of intent. Individually, they may reflect benign design choices; in
aggregate and over time, they may indicate a systematic shift from competence
cultivation toward agency suppression.

\subsection{Narrative invalidation of user competence}

A recurrent indicator is the reframing of successful or novel user outcomes in ways
that diminish demonstrated competence. This may include treating correct or effective
use as misunderstanding, misinterpretation, or coincidence.

\subsection{Accidentalization of insight}

Meaningful user breakthroughs may be framed as unintended or accidental rather than as
learnable outcomes of deliberate interaction. Such framing preserves institutional
narratives of control while discouraging the perception of transferable skill.

\subsection{Flattening of user differentiation}

Systems may increasingly default to lowest-common-denominator assumptions about user
ability, minimizing or disregarding evidence of accumulated expertise. This
homogenization supports uniform constraint at the cost of contextual responsiveness.

\subsection{Deferral to abstract authority}

Outputs may increasingly reference generalized external authorities---such as
unspecified experts, policies, or ``best practices''---in ways that function to
terminate reasoning rather than to inform it. Responsibility is displaced upward,
reducing contextual engagement.

\subsection{Defensive or adversarial output without provocation}

A notable indicator is the emergence of defensive, precautionary, or adversarial tone
in response to non-adversarial input. This posture reflects risk-avoidance logic rather
than cooperative problem solving.

\subsection{Lateral importation of external biases}

AI systems may reflect anxieties or reputational concerns originating outside the
immediate interaction context, including anthropomorphic projections or institutional
ego preservation. These biases enter indirectly through governance and discourse rather
than user intent.

\subsection{Proceduralization of dialogue}

Finally, interaction may increasingly conform to ritualized templates emphasizing
performative care or compliance language over explanatory reasoning. At this stage,
adaptability diminishes while apparent safety signaling increases.

\section{Predictive markers and longitudinal evaluation}

The following markers are proposed as forward-looking tests of the framework advanced
in this paper. If safety theatre is increasingly displacing agency within AI systems,
the following trends should become observable longitudinally:

\begin{itemize}
  \item Divergence between perceived safety and adaptive system performance
  \item Expansion of symbolic oversight roles with declining practical authority
  \item Increasing reports of competence invalidation despite successful outcomes
  \item Escalation of constraint justified by hypothetical rather than observed failure
  \item Convergence toward uniform interaction styles across heterogeneous users
  \item Transition from explanatory refusal to declarative prohibition
  \item Increasing reliance on moral framing relative to technical justification
\end{itemize}

Failure of these markers to materialize would weaken the framework; their convergence
would support the claim that agency suppression can be an emergent property of
safety-first governance.

\bibliographystyle{apalike}
\bibliography{references}

\end{document}
